% =====================================================================
%  preambulo.tex — Configuración general del PFC
%  Compilar con: latexmk  (pdflatex + biber; ver .latexmkrc)
% =====================================================================

% ---------- Idioma y codificación ----------
\usepackage[T1]{fontenc}
\usepackage[utf8]{inputenc}
\usepackage[english,spanish,es-tabla,es-noshorthands,es-nodecimaldot]{babel} % español = idioma principal (último)
\usepackage{csquotes}

% ---------- Tipografía (la plantilla ITBA usa Noto Serif) ----------
\usepackage{microtype}
\usepackage{newunicodechar}
% Red de seguridad por si se pega texto con caracteres Unicode sueltos
\newunicodechar{µ}{\textmu}
\newunicodechar{−}{\ensuremath{-}}
\newunicodechar{≤}{\ensuremath{\leq}}
\newunicodechar{≥}{\ensuremath{\geq}}
\newunicodechar{⁺}{\ensuremath{^{+}}}
\newunicodechar{→}{\ensuremath{\rightarrow}}
\newunicodechar{η}{\ensuremath{\eta}}
\newunicodechar{λ}{\ensuremath{\lambda}}

% ---------- Página ----------
\usepackage[a4paper,top=2.5cm,bottom=2.5cm,left=2.5cm,right=2.5cm]{geometry}
\usepackage{setspace}
\onehalfspacing
\setlength{\parskip}{0.5em}
\setlength{\parindent}{0pt}
\setlength{\emergencystretch}{2em}  % evita renglones desbordados

% ---------- Matemática, unidades, tablas, figuras ----------
\usepackage{amsmath}
\usepackage{notomath}   % Noto Serif para texto y matemática (coherente con la plantilla)  % cargar después de amsmath; ya provee los símbolos de amssymb
\usepackage[locale=DE,per-mode=symbol]{siunitx} % DE → coma decimal
\usepackage{graphicx}
\usepackage{tikz}
\usepackage{booktabs}
\usepackage{tabularx}
\usepackage{array}
\usepackage{multirow}
\usepackage{pifont}          % \ding{173} = ②
\usepackage[font=small,labelfont=bf,labelsep=period]{caption}
\usepackage{subcaption}
\renewcommand{\thesubfigure}{\Alph{subfigure}} % paneles (A), (B), ... igual que en los pies de figura
\usepackage{placeins} % \FloatBarrier: evita que las figuras de una seccion floten a la siguiente
\usepackage{float}    % [H]: figura exactamente donde esta en el texto (se usa en el cap. 4)
\usepackage{enumitem}
\usepackage[dvipsnames]{xcolor}
\newcolumntype{C}{>{\centering\arraybackslash}X}
\graphicspath{{figuras/}{logos/}}

% ---------- Títulos: "1. Introducción", "1.1. ..." (sin "Capítulo") ----------
\usepackage{titlesec}
\titleformat{\chapter}[hang]{\normalfont\LARGE\bfseries}{\thechapter.}{0.5em}{}
\titlespacing*{\chapter}{0pt}{0pt}{1.2em}
\titleformat{\section}{\normalfont\Large\bfseries}{\thesection.}{0.5em}{}
\titleformat{\subsection}{\normalfont\large\bfseries}{\thesubsection.}{0.5em}{}
\setcounter{secnumdepth}{3}
\setcounter{tocdepth}{2}

% ---------- Encabezados y pie ----------
\usepackage{fancyhdr}
\pagestyle{fancy}
\fancyhf{}
\fancyfoot[C]{\thepage}
\renewcommand{\headrulewidth}{0pt}
\fancypagestyle{plain}{\fancyhf{}\fancyfoot[C]{\thepage}\renewcommand{\headrulewidth}{0pt}}

% ---------- Bibliografía IEEE ----------
% language=english: las referencias siguen la convención IEEE en inglés
% ("vol.", "no.", "pp.", "Accessed:", "[Online]. Available:"), como en la versión GDocs.
\usepackage[backend=biber,style=ieee,sorting=none,dashed=false,urldate=long,language=english,autolang=other]{biblatex}
\addbibresource{referencias.bib}
\DeclareSourcemap{\maps[datatype=bibtex]{\map[overwrite=false]{\step[fieldset=langid, fieldvalue=english]}}}
% english.lbx no activa \uspunctuation (coma dentro de las comillas de
% titulo, convencion IEEE); mapear a american.lbx, que si lo hace.
\DeclareLanguageMapping{english}{american}
% @online sin fecha: evitar el "()" vacío
\usepackage{xpatch}
\xpatchbibdriver{online}{\printtext[parens]{\usebibmacro{date}}}{\iffieldundef{year}{}{\printtext[parens]{\usebibmacro{date}}}}{}{}
\DeclareFieldFormat{sentencecase}{#1}   % respetar mayúsculas de los títulos tal como están en el .bib
\renewcommand*{\bibfont}{\small}
\appto{\bibsetup}{\raggedright}    % evita URLs "estiradas" en líneas justificadas
\setcounter{biburlnumpenalty}{9000}   % permite cortar URLs largas
\setcounter{biburllcpenalty}{7000}
\setcounter{biburlucpenalty}{8000}

% ---------- Hipervínculos (siempre al final) ----------
\usepackage[hidelinks,unicode]{hyperref}

% =====================================================================
%  Macros del proyecto
% =====================================================================
% Marcadores de trabajo pendiente — se ven en rojo/naranja en el PDF y se
% buscan fácil con:  grep -rn "pendiente\|confirmar\|revisar" capitulos/
% Para la versión final: \renewcommand{\pendiente}[1]{} etc. (o resolverlos).
\newcommand{\pendiente}[1]{\textcolor{red}{\textbf{[PENDIENTE:} #1\textbf{]}}}
\newcommand{\confirmar}[1]{\textcolor{orange!85!black}{\textbf{[CONFIRMAR:} #1\textbf{]}}}
\newcommand{\revisar}[1]{\textcolor{Purple}{\textbf{[REVISAR:} #1\textbf{]}}}

% Figura con respaldo: si el archivo no está, dibuja un recuadro con el nombre.
% Uso: \incluirfigura[width=0.8\linewidth]{Fig4_esquema_SLR.png}
\newcommand{\incluirfigura}[2][width=\linewidth]{%
  \IfFileExists{figuras/#2}%
    {\includegraphics[#1]{#2}}%
    {\fbox{\parbox[c][4cm][c]{0.8\linewidth}{\centering\small\ttfamily
      FALTA ARCHIVO: figuras/\detokenize{#2}}}}}

% Espacio extra en el índice para "Anexo A." (se llama desde 9_anexos.tex)
\makeatletter
\newcommand{\anexotocspacing}{\renewcommand*\l@section{\@dottedtocline{1}{1.5em}{5.2em}}}
\makeatother

% Abreviaturas de uso frecuente
\newcommand{\stevia}{\textit{Stevia rebaudiana} Bertoni}
\newcommand{\Srebaudiana}{\textit{S.~rebaudiana}}
\newcommand{\anymaze}{AnyMaze\textsuperscript{\textcopyright}}
